% !TEX root = task-1-17-5.tex
% Editable task bank. Append new BankTask blocks at the end of the document.
\documentclass[12pt,a4paper]{article}
\usepackage{physics-textbook}
\usepackage{graphicx,needspace}
\usetikzlibrary{decorations.pathmorphing}
\newif\ifShowSolutions
\ifdefined\TutorEdition
  \ShowSolutionstrue
\else
  \ShowSolutionsfalse % Change to \ShowSolutionstrue to display solutions.
\fi
\newcounter{banktask}
\newsavebox{\BankTaskBox}
\newcommand{\BankTaskContent}[3]{%
  \noindent\textbf{\thebanktask.}\quad #2\par
  \ifShowSolutions
    \Needspace{50mm}% Keep the solution heading with its opening text or diagram.
    #3\par
    \ifdefined\TutorEdition
      \ifcsname TutorSource#1\endcsname
        \nopagebreak\smallskip{\small\csname TutorSource#1\endcsname\par}%
      \fi
    \fi
  \fi
}
\newcommand{\BankTask}[3]{%
  \refstepcounter{banktask}\par\medskip
  \sbox{\BankTaskBox}{\begin{minipage}{\linewidth}%
    \BankTaskContent{#1}{#2}{#3}\end{minipage}}%
  \ifdim\dimexpr\ht\BankTaskBox+\dp\BankTaskBox\relax>\dimexpr\textheight-20mm\relax
    \Needspace{35mm}\BankTaskContent{#1}{#2}{#3}%
  \else
    \Needspace{\dimexpr\ht\BankTaskBox+\dp\BankTaskBox\relax}%
    \noindent\usebox{\BankTaskBox}\par
  \fi
  \medskip
}
\begin{document}
\BookSubsection{1.17.5}{Свободные и вынужденные колебания}


% Source sdamgia-603: https://phys-ege.sdamgia.ru/problem?id=603
\BankTask{sdamgia-603}{
Груз колеблется на пружине, подвешенной вертикально к потолку, при этом максимальное расстояние от потолка до центра груза равно $H$, минимальное $h$. В точке, удаленной от потолка на расстояние $h$:

\noindent 1) кинетическая энергия шарика максимальна\par
\noindent 2) кинетическая энергия шарика минимальна\par
\noindent 3) потенциальная энергия пружины максимальна\par
\noindent 4) потенциальная энергия взаимодействия шарика с землей минимальна\par
}{
\smallskip\noindent\textbf{Решение}\par
В точке, удаленной от потолка на расстояние $h$, шарик меняет направление своего движения. Он перестает подниматься и начинает опускаться, поэтому скорость его в этой точке равна нулю, а значит, кинетическая энергия минимальна. Утверждения $3$ и $4$ относятся к положению шарика, когда он удален от потолка на расстояние $H$. В этот момент пружина максимально растянута, а шарик находится на минимальном расстоянии от земли.
Ответ: $2$.
}

% Source sdamgia-608: https://phys-ege.sdamgia.ru/problem?id=608
\BankTask{sdamgia-608}{
Шарик колеблется на пружине, подвешенной вертикально к потолку, при этом максимальное расстояние от потолка до центра шарика равно $H$, минимальное $h$. В точке, удаленной от потолка на расстояние $H$, максимальна:

\noindent 1) кинетическая энергия шарика\par
\noindent 2) потенциальная энергия пружины\par
\noindent 3) потенциальная энергия взаимодействия шарика с Землей\par
\noindent 4) сумма кинетической энергии шарика и взаимодействия шарика с Землей\par
}{
\smallskip\noindent\textbf{Решение}\par
В точке, удаленной от потолка на расстояние $H$, шарик меняет направление своего движения. Он перестает опускаться и начинает подниматься, поэтому скорость его в этой точке равна нулю, а значит, кинетическая энергия минимальна. В тоже время шарик находится на минимальном расстоянии от земли, следовательно, потенциальная энергия взаимодействия его с Землей также минимальна. Пружина, напротив, оказывается в этом положении максимально растянутой. Таким образом, верно утверждение $2$.
Ответ: $2$.
}

% Source sdamgia-8430: https://phys-ege.sdamgia.ru/problem?id=8430
\BankTask{sdamgia-8430}{
\begin{center}\begin{tikzpicture}[task graph]
\draw[task curve] (-1,2.5)--(1,2.5);\foreach \x in {-.9,-.65,...,.9}{\draw[task tick] (\x,2.5)--++(.2,.2);}
\draw[task curve] (0,2.5)--(-1.5,.35);\draw[task curve] (0,2.5)--(0,0);\draw[task curve] (0,2.5)--(1.5,.35);
\draw[task curve,fill=orange!25] (-1.5,.35) circle (.23);\draw[task curve,fill=orange!25] (0,0) circle (.23);\draw[task curve,fill=orange!55] (1.5,.35) circle (.23);
\end{tikzpicture}\end{center}Математический маятник с периодом колебаний $T$ отклонили на небольшой угол от положения равновесия и отпустили без начальной скорости (см. рис.). Через какое время (в долях периода) после этого кинетическая энергия маятника в первый раз достигнет минимума? Сопротивлением воздуха пренебречь.
}{
\smallskip\noindent\textbf{Решение}\par
Поскольку маятник отпустили с нулевой начальной скоростью, положение начального отклонения соответствует максимальному отклонению. За время, равное периоду, маятник успеет отклониться в противоположную сторону, после чего вернется в начальное положение. Минимальной кинетической энергии соответствует положение максимального отклонения. Впервые маятник окажется в нем через половину периода.

Ответ: $0{,}5$.
}

% Source sdamgia-601: https://phys-ege.sdamgia.ru/problem?id=601
\BankTask{sdamgia-601}{
Период колебаний потенциальной энергии пружинного маятника $1\,\text{с}$. Каким будет период ее колебаний, если массу груза маятника и жесткость пружины увеличить в $4$ раза?
}{
\smallskip\noindent\textbf{Решение}\par
Период колебаний потенциальной энергии пружинного маятника в два раза меньше периода колебаний самого маятника. В свою очередь, период колебаний пружинного маятника зависит только от отношения массы груза и жесткости пружины: 

$T=2 \pi \sqrt{\frac{m}{k}} .$

Таким образом, одновременное их увеличение в $4$ раза не приведет к изменению периода колебаний потенциальной энергии.

Ответ: $1$.
}

% Source sdamgia-526: https://phys-ege.sdamgia.ru/problem?id=526
\BankTask{sdamgia-526}{
Период колебаний потенциальной энергии горизонтального пружинного маятника $1\,\text{с}$. Каким будет период ее колебаний, если массу груза маятника увеличить в $2$ раза, а жесткость пружины вдвое уменьшить?
}{
\smallskip\noindent\textbf{Решение}\par
Период колебаний потенциальной энергии пружинного маятника пропорционален периоду колебаний груза, который определяется выражением:

$T=2 \pi \sqrt{\frac{m}{k}} .$

Следовательно, увеличение массы груза маятника в $2$ раза и уменьшение жесткости пружины в $2$ раза приведет к увеличению периода колебаний потенциальной энергии пружинного маятника в $2$ раза: $T_{\,\text{нов}}/T_{\,\text{стар}}= \sqrt{\frac{m'}{k'}} : \sqrt{\frac{m}{k}} = \sqrt{\frac{2m}{k/2}}: \sqrt{\frac{m}{k}} =2.$ Он окажется равным $T_{\,\text{нов}}=2T_{\,\text{стар}}=2\,\text{с}.$

Ответ: $2$.
}

% Source sdamgia-7618: https://phys-ege.sdamgia.ru/problem?id=7618
\BankTask{sdamgia-7618}{
Груз подвешен на легкой вертикальной пружине и совершает на ней колебания с частотой $\omega$ = $10\,\text{рад/с}$, двигаясь по вертикали. На какую длину растянется эта пружина, если аккуратно подвесить к ней тот же груз, не возбуждая колебаний? Ускорение свободного падения принять равным $10\,\text{м/с}^{2}$.
}{
\smallskip\noindent\textbf{Решение}\par
Жесткость пружины можем выразить как $k= \frac{mg}{L} ,$ где $L$  — удлинение пружины под действием груза.
Циклическая частота колебаний выражается как $\omega=\sqrt{\frac{k}{m}}=\sqrt{\frac{mg/L}{m}}=\sqrt{\frac gL}.$
Выразим отсюда удлинение: $L= \frac{g}{\omega ^2} =0{,}1\,\text{м}= 10\,\text{см}.$

Ответ: $10$.
}

% Source sdamgia-637: https://phys-ege.sdamgia.ru/problem?id=637
\BankTask{sdamgia-637}{
\begin{center}\begin{tikzpicture}[task graph,x=.65cm,y=.24cm]
\draw[task grid,xstep=1,ystep=2] (0,0) grid (7,14);\draw[task axis] (0,0)--(7.6,0) node[right] {$\nu$, Гц};\draw[task axis] (0,0)--(0,14.6) node[above] {$A$, см};
\draw[task curve] plot[smooth] coordinates {(0,1.1)(1,2)(1.6,4)(1.85,8)(2,10)(2.15,8)(2.4,4)(3,2)(4,1.1)(5,.7)(6,.5)};
\foreach \x in {2,4,6}{\node[below] at (\x,0) {$\x$};}\foreach \y in {2,4,6,8,10}{\node[left] at (0,\y) {$\y$};}\node[below left] at (0,0) {$0$};
\end{tikzpicture}\end{center}На рисунке изображена зависимость амплитуды установившихся колебаний маятника от частоты вынуждающей силы (резонансная кривая). Какова амплитуда колебаний этого маятника при резонансе?
}{
\smallskip\noindent\textbf{Решение}\par
Резонансом называется явление резкого возрастания амплитуды вынужденных колебаний при приближении частоты вынуждающей силы к собственной частоте маятника. Из графика видно, что резонанс происходит при значении частоты вынуждающей силы в $2\,\text{Гц}$, амплитуда колебаний маятника при этом равна $10\,\text{см}$.

Ответ: $10$.
}

% Source sdamgia-7850: https://phys-ege.sdamgia.ru/problem?id=7850
\BankTask{sdamgia-7850}{
\begin{center}\begin{tikzpicture}[task graph,x=1.4cm,y=.24cm]
\draw[task grid,xstep=.5,ystep=2] (0,0) grid (3.5,14);\draw[task axis] (0,0)--(3.8,0) node[right] {$\nu$, Гц};\draw[task axis] (0,0)--(0,14.6) node[above] {$A$, см};
\draw[task curve] plot[smooth] coordinates {(0,1.2)(.5,2)(.8,4)(.95,8)(1,10)(1.05,8)(1.2,4)(1.5,2)(2,1)(2.5,.6)(3,.5)};
\foreach \x/\s in {.5/{0{,}5},1/1,1.5/{1{,}5},2/2,2.5/{2{,}5},3/3}{\node[below] at (\x,0) {$\s$};}\foreach \y in {2,4,6,8,10}{\node[left] at (0,\y) {$\y$};}\node[below left] at (0,0) {$0$};
\end{tikzpicture}\end{center}На рисунке изображена зависимость амплитуды установившихся колебаний маятника от частоты вынуждающей силы (резонансная кривая). Частота вынуждающей силы вначале была равна $0{,}5\,\text{Гц}$, а затем стала равна $1{,}0\,\text{Гц}$.
Во сколько раз изменилась при этом амплитуда установившихся вынужденных колебаний маятника?
}{
\smallskip\noindent\textbf{Решение}\par
Как видно из графика амплитуда увеличилась с $2\,\text{см}$ до $10\,\text{см}$, то есть в $5$ раз.

Ответ: $5$.
}

% Source sdamgia-3131: https://phys-ege.sdamgia.ru/problem?id=3131
\BankTask{sdamgia-3131}{
Груз, подвешенный на пружине, совершает вынужденные гармонические колебания под действием силы, меняющейся с частотой $\nu .$ 

\TaskMatchingColumns{Физические величины}{А) Кинетическая энергия\\Б) Проекция скорости}{Частота их изменения}{1) $\frac1{2\nu}$\\2) $\nu$\\3) $2\nu$\\4) $\frac1{4\nu}$}
\begin{center}\begin{tabular}{|c|c|}\hline А&Б\\\hline &\\\hline\end{tabular}\end{center}
}{
\smallskip\noindent\textbf{Решение}\par
Под действием силы, меняющейся с частотой $\nu ,$ груз на пружине совершает вынужденные гармонические колебания с такой же частотой. Следовательно, закон изменения координаты груза со временем имеет вид $x=x_m \sin ( 2 \pi \nu t ) .$ Таким образом, закон изменения проекции скорости со временем: $V=2 \pi \nu x_m \cos ( 2 \pi \nu t ) =V_m \cos ( 2 \pi \nu t ) .$ Отсюда получаем, что частота изменения проекции скорости груза также равна $\nu$ (Б  — $2$). Кинетическая энергия груза изменяется по закону 

$E_{\,\text{кин}} ( t ) = \frac{mV ^2}{2} = \frac{mV_m ^2 \cos ^2 ( 2 \pi \nu t )}{2} = \frac{mV_m ^2}{2} \cdot \frac{1 + \cos ( 4 \pi \nu t )}{2} .$

 Следовательно, частота изменения кинетической энергии равна $2\nu$ (А  — $3$).

Ответ: $32$.
}

% Source sdamgia-35988: https://phys-ege.sdamgia.ru/problem?id=35988
\BankTask{sdamgia-35988}{
Пружинный маятник может совершать вынужденные колебания. В таблице приведена экспериментально полученная зависимость амплитуды $A$ установившихся вынужденных гармонических колебаний груза маятника от циклической частоты $\omega,$ на которой происходят эти колебания.

\begin{center}\begin{tabular}{|c|c|c|c|c|c|c|c|c|c|}\hline
$\omega,$ рад/с&$4$&$8$&$12$&$16$&$20$&$24$&$28$&$32$&$36$\\\hline
$A$, $\text{см}$&$2{,}1$&$2{,}4$&$2{,}9$&$4{,}4$&$11{,}6$&$10{,}5$&$3{,}3$&$1{,}9$&$1{,}2$\\\hline
\end{tabular}\end{center}

 Циклическая частота измеряется очень точно, абсолютная погрешность измерения амплитуды составляет $1\,\text{мм}$.

\noindent 1. При увеличении циклической частоты амплитуда установившихся вынужденных колебаний груза маятника все время возрастает.\par
\noindent 2. При всех значениях циклической частоты колебаний запас механической энергии системы одинаков.\par
\noindent 3. Резонансное значение амплитуды колебаний груза маятника наблюдается при величине циклической частоты, которая лежит между $16\,\text{рад/с}$ и $24\,\text{рад/с}$.\par
\noindent 4. При частоте $16\,\text{рад/с}$ максимальное значение модуля скорости груза маятника составляет приблизительно $70\,\text{см/с}$.\par
\noindent 5. При циклической частоте $36\,\text{рад/с}$ максимальное значение модуля скорости груза маятника меньше, чем при частоте $4\,\text{рад/с}$.\par
}{
\smallskip\noindent\textbf{Решение}\par
$1$.  Неверно. Циклическая частота постоянно возрастает, а амплитуда колебаний сначала возрастает, а затем убывает.
$2$.  Неверно. Так как полная механическая энергия маятника зависит соответственно формуле $E_{\,\text{полн}}= \frac{kA ^2}{2} ,$ она зависит прямо от квадрата амплитуды, которая в опыте меняется.
$3$.  Верно. Резонанс наблюдается при максимальной амплитуде, когда частота внешних воздействий совпадает с собственной частотой колебаний. Такие значения циклическая частота принимает в интервале от $16\,\text{рад/с}$ до $24\,\text{рад/с}$.
$4$.  Верно. Максимальное значение скорости маятника равно: 

$v _{\max}=\omega A=4{,}4 \cdot 16\approx 70 \,\text{см/с}.$

$5$.  Неверно. Максимальное значение скорости при циклической частоте $36\,\text{рад/с}$ равно: 

$v _{\max}=\omega A=1{,}2 \cdot 36\approx 43 \,\text{см/с}.$

При циклической частоте $4\,\text{рад/с}$: 

$v _{\max}=\omega A=2{,}1 \cdot 4\approx 8 \,\text{см/с}.$

Ответ: $34$.
}

\end{document}
